\section{From Blind Signatures to ARC}
\label{sec:arc-construction}

This section turns the blind signature of Section~\ref{sec:blind-signature} into an anonymous rate-limited credential
(ARC) system. We use the following terminology throughout:
\emph{issuance} refers to the blind signature protocol, while \emph{showing} refers to an ARC presentation.
The issuance proof is the \emph{pre-sign NIZK}; the showing proof is the \emph{post-sign NIZK}.

\subsection{From blind signatures to ARC}
A credential issued to a holder conceptually consists of:
\begin{itemize}[leftmargin=1.25em]
  \item the signed message \(\vecmu=\vecm\Vert \veck\), where \(\vecm\) are attributes and \(\veck\) is a uniformly sampled secret PRF key;
  \item the signature \(u\) satisfying \(A u = h_{\vecmu,(r_0,r_1)}(B)\) for some bounded randomness \((r_0,r_1)\).
\end{itemize}
The holder may additionally store local auxiliary values needed to prove the signature relation in zero knowledge. These
values are never revealed in showings.

\paragraph{Message structure.}
Every credential message has the abstract form
\[
  \vecmu = \vecm \Vert \veck.
\]
Here \(\vecm\) encodes the holder attributes and \(\veck\) is sampled uniformly and serves as the secret PRF key used for
rate limiting. At the proof-system layer this abstract message is represented by packed witness rows \((M,K)\), but the
credential syntax itself remains the quotient-ring statement of Section~\ref{sec:blind-signature}.

\subsection{Rate limiting via the secret key \(\veck\)}
Rate limiting is implemented by deterministic tags derived from a PRF keyed by the credential's secret \(\veck\). For a
public nonce \(\nonce\in\mathcal{D}\), the holder computes
\[
  \prftag = \PRF(\veck,\nonce).
\]
Intuitively:
\begin{itemize}[leftmargin=1.25em]
  \item if the holder never repeats a nonce, then \(\prftag\) is pseudorandom and yields unlinkable showings;
  \item if the holder repeats a nonce, the tag repeats deterministically, enabling the verifier to detect double-spending
  for that nonce.
\end{itemize}
The verifier enforces its rate policy by restricting acceptable nonces and maintaining state that records previously
accepted \((\nonce,\prftag)\) pairs for a given policy context.

\subsection{Concrete PRF instantiation (Poseidon2-like)}
\label{subsec:poseidon-prf}

\providecommand{\Tr}{\mathrm{Tr}}

We instantiate the PRF from Definition~\ref{def:prf} using the Poseidon2-like construction of
Definitions~\ref{def:poseidon-prf} and \ref{def:poseidon-perm}. Algorithm~\ref{alg:prf} summarizes the computation.

\begin{algorithm}[H]
\caption{$\PRF(k,n)$ (Poseidon2-like with feed-forward + truncation)}
\label{alg:prf}
\begin{algorithmic}[1]
\State $x \gets k\Vert n$; $x^{(0)}\gets x$
\For{$i=1$ to $R_F/2$} \State $x\gets E_i(x)$ \EndFor
\For{$i=1$ to $R_P$} \State $x\gets I_i(x)$ \EndFor
\For{$i=R_F/2+1$ to $R_F$} \State $x\gets E_i(x)$ \EndFor
\State $y \gets x + x^{(0)}$ \Comment{feed-forward}
\State \Return $\Tr(y)$ \Comment{first $\ell_{\prftag}$ lanes}
\end{algorithmic}
\end{algorithm}

\paragraph{Proof-friendly arithmetization.}
In the post-sign NIZK we enforce $\prftag=\PRF(\veck,\nonce)$ using a checkpointed arithmetization: we commit
selected S-box outputs and let the verifier replay all linear wiring between checkpoints from openings and public
constants. Section~\ref{subsec:prf-checkpointed} states the constraint families; Appendix~\ref{app:prf} records the
checkpoint schedule and extended notes.

\begin{comment}
Preserved from Carsten's 10-04 source: this entire PRF-instantiation subsection was wrapped in a comment environment at
that revision boundary. The active paper keeps the subsection live because the later rewrite uses it to anchor the
SmallWood PRF constraint discussion.
\end{comment}

\subsection{Showing protocol}
\label{subsec:showing-protocol}
\label{subsec:arc-protocol}

\begin{algorithm}[H]
\caption{\(\Show\): ARC presentation (canonical protocol)}
\label{alg:showing}
\begin{algorithmic}[1]
\Require Holder has credential \((\vecm,\veck,r_0,r_1,u)\). Public parameters include \(A,B,\BoundMsg\) and the PRF description.
\Require Input nonce \(\nonce\in\mathcal{D}\) that has not been used before for this credential.
\Ensure Presentation \((\prftag,\nonce,\pi_{\mathsf{post}})\).
\State Compute \(\prftag\gets \PRF(\veck,\nonce)\).
\State Produce a \emph{post-sign} NIZK \(\pi_{\mathsf{post}}\) proving knowledge of \((u,\vecm,\veck,r_0,r_1)\) such that:
\begin{enumerate}[leftmargin=1.6em]
  \item \(A u = h_{\vecmu,(r_0,r_1)}(B)\) for \(\vecmu=\vecm\Vert \veck\);
  \item \(\prftag = \PRF(\veck,\nonce)\);
  \item \(\vecm,\veck,r_0,r_1\) are coefficient-wise in \([-\BoundMsg,\BoundMsg]\) and \(u\) satisfies the public shortness bound.
\end{enumerate}
\State Output \((\prftag,\nonce,\pi_{\mathsf{post}})\).
\end{algorithmic}
\end{algorithm}
\cb{FIX THIS}

At the ARC level, the intended abstract witness remains 
\((u,\vecm,\veck,r_0,r_1)\), together with
\(\vecmu=\vecm\Vert \veck\), satisfying the signature equation and
 the PRF-tag equation. Section~\ref{sec:smallwood-model} later specifies the compiled SmallWood relation used to realize
that abstract witness statement. The coefficient-domain statement used in the security discussion below is recovered only
through the theorem suite of Section~\ref{subsec:statement-strength}; we do not use any stronger identification between
the compiled replay-space statement and the abstract showing relation.

The verifier runs \(\Verify\) by checking \(\pi_{\mathsf{post}}\) and consulting a policy state that tracks used
nonces and tags. The verifier rejects if the proof fails or if the nonce-tag pair violates the local policy.
\cb{No message $\mu$ here}

\begin{comment}
\subsubsection{ARC protocol}
\label{subsec:arc-protocol}
We assume the same setup as in Section \ref{subsec:bs-protocol}. Moreover, define the relations
\[
\relIssueARC,\qquad \relShow.
\]
In the 10-04 source, \(\relIssueARC\) explicitly committed \((\vecmu,\veck,\vecr_0^{\Holder},\vecr_1^{\Holder},\bar r)\)
and required
\[
\vect = h_{\vecmu \Vert \veck,(\vecr_0,\vecr_1)}(B),
\]
while \(\relShow\) explicitly bound
\[
\prftag = \PRF(\veck,\nonce),\qquad
\mA \vecu = h_{\vecmu\Vert \veck ,(\vecr_0,\vecr_1)}(B).
\]
\cb{FIX THIS}
That version then presented explicit \(\Issue\), \(\Show\), and \(\Verify\) routines in terms of
\(\relIssueARC,\relShow\), prior to the later theorem-clean rewrite that compiles showing through the Section~5
SmallWood relation instead.

\subsection{Showing protocol}
The 10-04 source also stated an older abstract showing algorithm in which \(\pi_{\mathsf{post}}\) directly proved
knowledge of \((u,\mu,\veck,r_0,r_1)\) satisfying the signature/hash relation, the PRF-tag relation, and the public
bounds, rather than referring forward to the compiled Section~5 relation.
\end{comment}

\subsection{PRF role and rationale}
The PRF serves two simultaneous roles:
\begin{itemize}[leftmargin=1.25em]
  \item \textbf{Public rate-limiting handle.} The tag is deterministic in \((\veck,\nonce)\) so the verifier can
  enforce stateful policies without interacting with the issuer and without learning the credential message.
  \item \textbf{Unlinkability across nonces.} For fresh nonces, tags are pseudorandom even given previous tags, so a
  verifier cannot link two presentations to the same credential unless the policy context forces reuse.
\end{itemize}
Crucially, \(\veck\) is part of the signed message, so the holder cannot change the PRF key between showings without
producing a new valid signature witness. The post-sign NIZK binds the public tag to that signed hidden key while keeping
the key itself hidden. This is the mechanism by which verifier-side state can enforce rate policies through the public
pair \((\nonce,\prftag)\).

\subsection{ARC correctness and conditional security decomposition}

\paragraph{Honest correctness.}
If issuance completes, the holder stores a witness \((u,\vecm,\veck,r_0,r_1)\), together with
\(\vecmu=\vecm\Vert \veck\), satisfying the abstract signature relation
of Section~\ref{sec:blind-signature}. For any public nonce \(\nonce\), the holder computes
\(\prftag = \PRF(\veck,\nonce)\) and derives the corresponding witness for the compiled post-sign relation of
Section~\ref{sec:smallwood-model}. By completeness of the post-sign NIZK, the verifier accepts any honest presentation
for which the local policy check on \((\nonce,\prftag)\) also accepts.

\paragraph{Cryptographic decomposition of the ARC layer.}
The ARC layer combines five ingredients:
\begin{enumerate}[leftmargin=1.25em]
    \item the non-blind theorem-backed signature layer of Lemma~\ref{lem:vsis-unforgeability};
    \item blindness and one-more unforgeability of the blind issuance protocol, which remain to be reduced;
    \item straightline-extractable knowledge soundness and zero knowledge of the post-sign NIZK for the compiled relation of Section~\ref{sec:smallwood-model};
    \item pseudorandomness and correctness of the PRF; and
    \item verifier-side policy state recording previously accepted \((\nonce,\prftag)\) pairs.
\end{enumerate}
Only the first and third ingredients are theorem-backed in the present draft. The fourth is an external assumption on
the chosen PRF instantiation, the second remains a reduction goal, and the fifth is an application-layer enforcement
rule rather than a stand-alone cryptographic theorem.

\paragraph{Conditional policy soundness.}
Knowledge soundness of the post-sign proof yields a witness for the compiled showing relation of
Section~\ref{sec:smallwood-model}. Section~\ref{subsec:statement-strength} is the sole bridge from that compiled
replay-space statement to the coefficient-domain statement used here. In particular,
Corollary~\ref{cor:showing-proof-consequence} yields a witness \((u,T,M,K,R_0,R_1,Z)\) satisfying
\[
B_3Au-(Au)R_1-B_0-B_1(M+K)-B_2R_0=0
\]
in \(R_q\), together with
\[
(B_3-R_1)\Had Z = 1,
\qquad
Au = h_{M+K,(R_0,R_1)}(B).
\]
Under that signature-side condition, together with one-more unforgeability of the blind issuance layer and correctness
of the PRF, the signed key \(\veck\) determines the public tag uniquely for each nonce. Consequently, a verifier that
records spent \((\nonce,\prftag)\) pairs rejects same-nonce replays deterministically.

\paragraph{Conditional presentation unlinkability.}
Presentation unlinkability in Definition~\ref{def:pres-unlink} is not implied by PRF security and post-sign zero
knowledge alone. In that experiment the adversary controls the issuer during the two issuance sessions. Accordingly,
unlinkability also depends on an issuance-layer privacy property, namely blindness or a comparably strong issuer-view
hiding statement. Under that issuance-layer privacy assumption, post-sign zero knowledge hides the witness used during
showing, and PRF pseudorandomness makes the public tags computationally random across distinct fresh nonces. This yields
the intended unlinkability claim for showing transcripts, subject to the explicit nonce-reuse caveat.

\paragraph{Attribute privacy and selective disclosure.}
The message part \(\vecm\) can encode attributes; the post-sign NIZK can be extended to prove predicates about \(\vecm\) and
to selectively reveal components, without changing the issuance protocol. We focus on the core rate-limiting
construction and defer predicate engineering to the SmallWood programming model in Section~\ref{sec:smallwood-model}.

\begin{comment}
\subsection{ARC security sketch}
\paragraph{Correctness.}
If issuance completes and the holder uses a fresh nonce, then the post-sign statement is satisfiable by construction.

\paragraph{Unforgeability and policy soundness.}
By knowledge soundness, any accepting presentation yields a witness satisfying both the signature/hash relation and the
PRF-tag relation. Under one-more unforgeability of the blind signature, an adversary cannot produce accepting showings
for credentials it did not obtain via issuance. Policy soundness is enforced by verifier state because
\(\prftag=\PRF(\veck,\nonce)\) is uniquely determined by \((\veck,\nonce)\).

\paragraph{Unlinkability.}
Presentation unlinkability is formalized in Definition~\ref{def:pres-unlink}. For distinct fresh nonces, PRF security
and zero knowledge hide the credential secret, except for the intended nonce-reuse caveat.

\paragraph{Constraint view (what the post-sign NIZK must prove).}
To prove \(\mathsf{tag}=\PRF(k,\mathsf{nonce})\) inside SmallWood, where \(k\) is the signed secret part of the
message, we arithmetize the S-box constraints, the public linear layers, and the feed-forward plus truncation
constraints binding the public tag to the final state.
\end{comment}
